% Options for packages loaded elsewhere
\PassOptionsToPackage{unicode}{hyperref}
\PassOptionsToPackage{hyphens}{url}
\documentclass[
]{article}
\usepackage{xcolor}
\usepackage{amsmath,amssymb}
\setcounter{secnumdepth}{-\maxdimen} % remove section numbering
\usepackage{iftex}
\ifPDFTeX
  \usepackage[T1]{fontenc}
  \usepackage[utf8]{inputenc}
  \usepackage{textcomp} % provide euro and other symbols
\else % if luatex or xetex
  \usepackage{unicode-math} % this also loads fontspec
  \defaultfontfeatures{Scale=MatchLowercase}
  \defaultfontfeatures[\rmfamily]{Ligatures=TeX,Scale=1}
\fi
\usepackage{lmodern}
\ifPDFTeX\else
  % xetex/luatex font selection
\fi
% Use upquote if available, for straight quotes in verbatim environments
\IfFileExists{upquote.sty}{\usepackage{upquote}}{}
\IfFileExists{microtype.sty}{% use microtype if available
  \usepackage[]{microtype}
  \UseMicrotypeSet[protrusion]{basicmath} % disable protrusion for tt fonts
}{}
\makeatletter
\@ifundefined{KOMAClassName}{% if non-KOMA class
  \IfFileExists{parskip.sty}{%
    \usepackage{parskip}
  }{% else
    \setlength{\parindent}{0pt}
    \setlength{\parskip}{6pt plus 2pt minus 1pt}}
}{% if KOMA class
  \KOMAoptions{parskip=half}}
\makeatother
\usepackage{longtable,booktabs,array}
\newcounter{none} % for unnumbered tables
\usepackage{calc} % for calculating minipage widths
% Correct order of tables after \paragraph or \subparagraph
\usepackage{etoolbox}
\makeatletter
\patchcmd\longtable{\par}{\if@noskipsec\mbox{}\fi\par}{}{}
\makeatother
% Allow footnotes in longtable head/foot
\IfFileExists{footnotehyper.sty}{\usepackage{footnotehyper}}{\usepackage{footnote}}
\makesavenoteenv{longtable}
\setlength{\emergencystretch}{3em} % prevent overfull lines
\providecommand{\tightlist}{%
  \setlength{\itemsep}{0pt}\setlength{\parskip}{0pt}}
\usepackage{bookmark}
\IfFileExists{xurl.sty}{\usepackage{xurl}}{} % add URL line breaks if available
\urlstyle{same}
\hypersetup{
  hidelinks,
  pdfcreator={LaTeX via pandoc}}

\author{}
\date{}

\begin{document}

\section{MPTA-Repair: Multi-Property Clock-Bound Repair for
Control-Domain Timed
Automata}\label{mpta-repair-multi-property-clock-bound-repair-for-control-domain-timed-automata}

\subsection{Abstract}\label{abstract}

Timed automata are widely used to model and verify real-time control
systems, where timing-safety properties such as bounded response,
deadline satisfaction, timeout handling, and unsafe-state avoidance are
essential to system reliability. In practical industrial settings,
however, control models are rarely simple single-property artifacts.
They often consist of multiple interacting components, clocks,
synchronizations, guards, invariants, and observer properties, and the
model must satisfy a set of timing requirements simultaneously. When
such a model violates several properties, engineers may need to analyze
multiple timed diagnostic traces and determine which clock-bound
constants should be repaired without breaking other requirements or
changing the intended functional behavior. Existing trace-based repair
workflows provide valuable support for local timed-automata repair, but
they are not specifically designed for the joint repair of models in
which multiple properties, multiple counterexamples, and multiple repair
variables interact.

This paper presents MPTA-Repair, a counterexample-guided workflow for
multi-property clock-bound repair of control-domain timed automata.
MPTA-Repair modifies numerical clock bounds in transition guards and
location invariants, and uses MaxSMT-based repair generation to
synthesize small candidate edits. To improve repair efficiency and
robustness, it jointly considers relations among violated properties,
diagnostic traces, and repair variables during iterative refinement.
Each candidate repair is validated by full-property regression
verification, and a TARTAR-style admissibility check based on functional
equivalence is used to ensure that the repaired model preserves the
intended untimed behavior. We construct a benchmark suite by collecting
control-domain timed-automata models from railway, automotive, cardiac
pacing, autonomous-driving, communication, ventilation, and aerospace
scheduling scenarios, and by injecting clock-bound faults to generate
faulty model instances. Experimental results show that MPTA-Repair
achieves better repair effectiveness than an iterative TARTAR-style
baseline, especially on control-domain models where local
single-property repairs frequently fail global validation.

\textbf{Keywords---} Timed automata; Multi-property repair; Clock-bound
repair; Real-time control systems.

\subsection{I. Introduction}\label{i.-introduction}

Real-time control software is increasingly deployed in domains where
timing faults may directly compromise safety, availability, or mission
success. In railway systems, timed automata have been used to model
train-gate controllers, communication protocols, and railway safety
layers {[}24{]}. In automotive and autonomous-driving systems, they have
been used to reason about gear-control logic, road-junction behavior,
and time-bounded decision rules {[}21{]}. In medical cyber-physical
systems, timed models have been used to analyze pacemakers and
ventilator controllers, where physiological timing windows and actuation
deadlines are safety-critical {[}22{]}. In aerospace and avionics,
timed-automata models have supported schedulability and deadline
analysis for control software and distributed real-time architectures
{[}27{]}. Across these domains, correctness depends not only on the
order of discrete events, but also on when these events occur. A
controller that eventually raises an alarm, applies a pulse, sends a
message, or releases a resource may still be incorrect if the action
occurs too early, too late, or for the wrong duration.

Timed automata provide a well-established modeling formalism for this
class of systems {[}1{]}. Real-time systems are often represented as
networks of timed automata, where clocks, guards, invariants,
synchronizations, and resets describe both discrete behavior and
quantitative timing constraints {[}2{]}. Verification is then performed
against temporal properties that may include safety, bounded response,
deadline, error-state reachability, alarm-duration, and deadlock-freedom
requirements. When a property is violated, a real-time model checker
such as UPPAAL, Kronos, or opaal can produce a timed diagnostic trace
(TDT), i.e., a timed counterexample that shows how the model reaches a
violating state {[}5{]}.

However, diagnosis is only the first step. Interpreting a TDT and
turning it into a correct model edit remains a demanding engineering
task. Realistic control models often contain many locations,
transitions, clocks, guards, invariants, synchronizations, and resets. A
TDT exposes one feasible execution that violates a property, but it does
not directly identify which clock constraint is wrong, whether the
defect lies in a guard or an invariant, or how a bound should be changed
to restore the intended timing behavior {[}6{]}. This ambiguity becomes
more pronounced when the model is checked against multiple properties.
Control-system specifications are usually not single assertions, but
collections of timing requirements that jointly describe safety,
response, scheduling, and reliability obligations {[}29{]}. A local edit
that eliminates one diagnostic trace may still leave related traces
feasible, or may even make another property false. Thus, automated
repair for control models must go beyond explaining a single
counterexample and must validate each candidate repair against the
complete property set.

Existing timed-automata repair techniques provide an important
foundation. Clock-bound repair for timed systems encodes a TDT into
linear real arithmetic, introduces variation variables for clock bounds,
and uses MaxSMT solving to compute small modifications to guards and
invariants {[}6{]}. TarTar extends this idea into an automatic repair
analysis tool that suggests syntactic repairs and checks whether the
repaired model preserves the untimed functional behavior of the original
network timed automaton {[}6{]}. Other related directions include
parameter synthesis for parametric timed automata {[}8{]}, clock-guard
repair through abstraction and testing {[}10{]}, robustness analysis for
uncertain timing constants {[}11{]}, and SAT/SMT-based model or program
repair {[}15{]}. These methods show that automated repair of timed
models is feasible, but most existing TDT-based repair workflows remain
centered on one violated property or one diagnostic trace at a time.

The engineering challenge addressed in this paper is different but
complementary to existing single-trace repair and synthesis techniques
{[}10{]}. A straightforward repair loop can process one violated
property, collect one TDT, compute a local repair, instantiate a
candidate model, and then run regression verification. This loop is
useful as a baseline, but it may become inefficient or unstable when one
timing fault affects several properties or several related traces.
Similar traces may repeatedly encode overlapping parts of the same
faulty behavior; a candidate that eliminates one trace may still leave
other violations feasible; and a bound changed for one property may
affect another property when the same clock, transition, location, or
synchronization participates in both executions. If all editable bounds
are considered equally plausible repair targets, the solver may explore
many weakly related candidates, including edits that are syntactically
valid but distant from the timing decisions exposed by the observed
violations. Such repairs may pass a local check while giving engineers
little insight into the shared cause of the failures.

This paper presents \textbf{MPTA-Repair}, an automatic repair workflow
for multi-property, multi-counterexample clock-bound repair of
timed-automata control models. Given a faulty model, a set of timing
properties, and diagnostic traces produced by model checking,
MPTA-Repair suggests small modifications to numerical clock bounds in
transition guards and location invariants. The repair space is
intentionally restricted to clock-bound constants, rather than arbitrary
structural edits, specification changes, synchronization changes, or
reset changes. The workflow uses lightweight relation analysis among
properties, diagnostic traces, and editable clock bounds to guide repair
search, but this analysis is only an optimization strategy. Every
accepted candidate must satisfy the complete property set through
regression verification and must pass a TARTAR-style admissibility check
based on functional equivalence. In this way, MPTA-Repair aims to
generate repairs that are not only local counterexample fixes, but also
system-level timing corrections validated against the full set of
control requirements.

We evaluate MPTA-Repair on a suite of control-domain timed-automata
models covering railway alarm control {[}25{]}, train-gate control
{[}4{]}, automotive gear control, dual-chamber pacemaker control,
autonomous road-junction rules, railway communication timeout handling,
mechanical ventilator control, and aerospace task scheduling. The
evaluation measures repair effectiveness, repair size, runtime,
regression results, and behavior preservation. We compare MPTA-Repair
with an iterative TARTAR-style baseline under the same full-property
regression and admissibility requirements, and report practical
observations from applying automated repair to industrial-style control
models.

This paper makes the following contributions:

\begin{enumerate}
\def\labelenumi{\arabic{enumi}.}
\tightlist
\item
  We formulate clock-bound repair for timed-automata control models as a
  multi-property and multi-counterexample repair problem.
\item
  We present MPTA-Repair, an automatic repair workflow that uses
  diagnostic traces and lightweight relation analysis to guide
  clock-bound repair.
\item
  We implement a prototype repair pipeline combining diagnostic trace
  analysis, MaxSMT-based clock-bound repair, full-property regression
  verification, and TARTAR-style admissibility checking.
\item
  We evaluate the approach on a diverse control-domain benchmark suite
  and compare it with an iterative single-property repair baseline under
  the same acceptance criteria.
\item
  We report practical lessons on applying automated timed-automata
  repair to control models, especially regarding property selection,
  repair interpretability, and the need for full regression after each
  candidate repair.
\end{enumerate}

The remainder of this paper is organized as follows. Section II provides
the background and problem setting, including timed automata, diagnostic
traces, clock-bound repair, and the multi-property repair objective.
Section III presents MPTA-Repair and its relation-guided repair
workflow. Section IV describes the prototype implementation and
evaluation setup. Section V reports the evaluation results and practical
lessons learned from the control-domain models. Section VI discusses
threats to validity and related work. Section VII concludes the paper.

\subsection{II. Background and Problem
Setting}\label{ii.-background-and-problem-setting}

\subsubsection{A. Timed Automata and
Networks}\label{a.-timed-automata-and-networks}

We use timed automata as the formal model for real-time control behavior
{[}1{]}. Let \texttt{C} be a finite set of clocks. A clock constraint
over \texttt{C} is a conjunction of atomic constraints of the form
\texttt{c\ ∼\ n}, where \texttt{c\ ∈\ C},
\texttt{∼\ ∈\ \{\textless{},\ ≤,\ =,\ ≥,\ \textgreater{}\}}, and
\texttt{n} is a non-negative numerical bound {[}6{]}. We write
\texttt{B(C)} for the set of such clock constraints. A timed automaton
is a tuple \texttt{T\ =\ (L,\ l0,\ C,\ Σ,\ Θ,\ I)}, where \texttt{L} is
a finite set of locations, \texttt{l0\ ∈\ L} is the initial location,
\texttt{Σ} is a set of action labels, \texttt{Θ} is a finite set of
transitions, and \texttt{I\ :\ L\ →\ B(C)} assigns an invariant to each
location {[}2{]}. A transition \texttt{θ\ ∈\ Θ} can be written as
\texttt{(l,\ g,\ a,\ r,\ l\textquotesingle{})}, where \texttt{l} and
\texttt{l\textquotesingle{}} are the source and target locations,
\texttt{g\ ∈\ B(C)} is the guard, \texttt{a\ ∈\ Σ} is the action label,
and \texttt{r\ ⊆\ C} is the set of clocks reset by the transition.

The semantics is defined over states \texttt{(l,\ u)}, where \texttt{l}
is a location and \texttt{u\ :\ C\ →\ R≥0} is a clock valuation. A delay
transition lets time elapse in the current location while its invariant
remains true. An action transition can be taken when its guard is
satisfied; it moves the automaton to a target location and resets the
specified clocks. Thus, guards constrain when discrete actions are
enabled, while invariants constrain how long the automaton can remain in
a location. Networks of timed automata compose several such automata,
often representing different parts of a control system, such as a
controller, plant, environment, communication protocol, monitor, or
scheduler {[}3{]}.

Timed automata are particularly suitable for control-domain models
because many requirements are expressed as numerical timing thresholds.
For instance, vehicle controllers may need to complete a gear-shifting
sequence within a bounded response time {[}21{]}; railway controllers
may need to satisfy gate, alarm, or communication timeout requirements
{[}24{]}; medical devices may encode physiological timing windows and
refractory intervals {[}22{]}; and avionics or aerospace scheduling
models may encode execution-time bounds and deadlines {[}28{]}. In these
models, a numerical bound in a guard or invariant often represents a
design threshold, timeout value, sampling period, alarm duration, or
resource deadline.

This paper focuses on faults in such numerical clock bounds. The
assumption is not that model structure is irrelevant, but that in many
engineering settings the designer has already captured the intended
control structure while some timing constants may be inaccurate due to
calibration, dimensioning, transcription, or design-space uncertainty.
This fault model is consistent with the motivation of clock-bound
repair: a clock-bound change can correct a clerical modeling mistake or
indicate that the timing resources of the design need to be
re-dimensioned {[}6{]}. It is also consistent with the broader intuition
behind mutation-based repair and testing that many defects can be
modeled as small deviations from an otherwise nearly correct artifact
{[}32{]}.

\subsubsection{B. Timed Diagnostic
Traces}\label{b.-timed-diagnostic-traces}

Timed-automata models are checked against timing properties using
real-time model checkers {[}3{]}. In this paper, we focus on timed
safety properties, i.e., properties that require all reachable states to
satisfy a specified invariant condition. Such properties commonly
constrain whether a location may be reached, whether a clock may exceed
a bound in a location, or whether an error state is reachable. More
complex requirements can also be represented using observer automata or
monitor templates that reduce them to safety-style checks {[}29{]}.

A symbolic timed trace is a finite sequence of transitions
\texttt{S\ =\ θ0,\ ...,\ θn−1} {[}6{]}. A realization of \texttt{S} is a
sequence of non-negative delays that makes the transition sequence
executable from the initial state while respecting guards, resets, and
invariants. A trace is feasible if it has at least one realization. If a
feasible trace has a realization whose final state violates a timed
safety property \texttt{Π}, then the trace is a timed diagnostic trace
for \texttt{Π} {[}6{]}.

A TDT is valuable because it gives concrete evidence of a timing
violation. It records a sequence of timed and discrete behavior leading
to a property-violating state, depending on the trace representation
produced by the model checker. However, it is not a root-cause
explanation {[}6{]}. A TDT does not state which numerical bound is
wrong, whether the relevant fault is in a guard or invariant, or how the
model should be modified. Multiple clock constraints can jointly enable
the violating realization, and one incorrect bound can contribute to
more than one TDT. Therefore, TDTs are best understood as repair
evidence rather than complete repair instructions.

This distinction is important in control models. A repair that simply
blocks the discrete path of a TDT may remove the observed violation but
also remove an intended functional behavior, such as a control action,
communication response, scheduling transition, or medical actuation. A
valid repair should instead correct the timing behavior while preserving
the relevant untimed functionality of the model {[}6{]}.

\subsubsection{C. Clock-Bound Repair}\label{c.-clock-bound-repair}

Clock-bound repair modifies numerical bounds in clock constraints while
preserving the surrounding automaton structure {[}6{]}. In this paper, a
repairable bound is an occurrence of a numerical constant \texttt{n} in
an atomic clock constraint \texttt{c\ ∼\ n} that appears in a transition
guard or a location invariant. A repair changes \texttt{n} to a new
value \texttt{n\textquotesingle{}}, or equivalently applies a variation
\texttt{v} so that the repaired bound becomes \texttt{n\ +\ v}.

The current MPTA-Repair pipeline does not modify property formulas, add
or delete locations, add or delete transitions, change synchronization
labels, replace clock references, or add/remove clock resets {[}10{]}.
This scope is deliberate. We aim to repair timing-parameter faults in a
model whose discrete control structure is treated as the intended
design. This avoids repairs that satisfy a property by changing the
specification, altering the modeled protocol structure, or disabling
intended behavior {[}6{]}. It also keeps the repair output
understandable to engineers, because each edit corresponds to a concrete
timing threshold in the model.

TDT-based clock-bound repair encodes the feasibility of a diagnostic
trace as constraints over delays and clock values {[}6{]}. Bound
variation variables are then introduced for clock constraints occurring
in guards and invariants. MaxSMT solving is used to find a small set of
bound changes that prevents the trace from having a realization that
violates the property {[}13{]}. In the original single-trace setting,
this yields a local clock-bound repair: the same discrete trace remains
executable, but no timing realization of that trace witnesses the
considered property violation.

A local repair is not automatically valid for the full model. Changing
clock bounds can affect other executions, other properties, and the
reachability of functional behavior. The TARTAR approach therefore
defines admissibility using functional equivalence: the repaired model
should preserve the untimed language of the original model, so that no
action traces are added or removed by the repair {[}6{]}. MPTA-Repair
adopts this admissibility criterion as the behavior-preservation gate
for accepted candidates.

\subsubsection{D. Multi-Property Repair
Objective}\label{d.-multi-property-repair-objective}

Control systems are usually validated by a set of timing properties
rather than by a single assertion {[}29{]}. These properties jointly
express the expected safety and reliability envelope of the model, such
as bounded response, absence of unsafe states, timeout handling,
deadline satisfaction, and deadlock freedom {[}34{]}. Let
\texttt{M\_bad} be a faulty timed-automata control model, and let
\texttt{Φ\ =\ \{φ1,\ ...,\ φm\}} be the property set used to validate
it. The model may violate one or more properties, and each violated
property may produce one or more TDTs. The goal is to construct a
repaired model \texttt{M\_rep} by modifying a small number of
clock-bound constants in guards and invariants.

A candidate repair is accepted only if it satisfies three requirements.
First, the repaired model must remain syntactically valid as a
timed-automata model; in particular, repaired clock bounds must remain
non-negative, and rational bounds introduced by solving must be
normalized if required by the target tool representation. Second, the
repaired model must satisfy the complete property set \texttt{Φ} under
regression verification, not only the property that produced the current
TDT. Third, the repaired model must preserve the functional behavior of
the original model using the TARTAR-style admissibility test based on
functional equivalence {[}6{]}.

The multi-property setting imposes obligations beyond local single-trace
repair {[}35{]}. Different properties may share clocks, locations,
synchronization actions, or execution fragments. A repair computed for
one property can therefore affect the behavior relevant to another
property. Similarly, multiple TDTs may reflect the same underlying
timing defect, so repairing them independently may duplicate solving
effort or produce inconsistent candidate edits. MPTA-Repair does not
treat these properties and traces as unrelated subproblems. Instead, it
analyzes them as a joint repair task, generates candidates using
multi-property and multi-counterexample evidence, and accepts a repair
only after iterative regression verification and admissibility checking.

\subsection{III. MPTA-Repair}\label{iii.-mpta-repair}

\subsubsection{A. Overview}\label{a.-overview}

MPTA-Repair is a counterexample-guided repair workflow for
multi-property, multi-counterexample clock-bound repair of
timed-automata control models {[}19{]}. Its input is a faulty model
\texttt{M}, a property set \texttt{Φ\ =\ \{φ1,\ ...,\ φm\}}, and
configuration parameters such as the maximum number of refinement
rounds, candidate limits, and a global timeout. The repairable bounds
are not given as a separate user input; they are extracted from the
numerical clock bounds occurring in transition guards and location
invariants {[}6{]}. The output is either a repaired model
\texttt{M\textquotesingle{}} together with a set of clock-bound edits,
or a failure report explaining why no admissible full-property repair
was found within the configured budget.

The workflow is designed around two principles. First, candidate repair
generation may be driven by local diagnostic traces, but candidate
acceptance must be global. A candidate that repairs one violated
property may cause another property to fail, because the same clock,
location, transition, or synchronization may participate in several
timing requirements. MPTA-Repair therefore keeps collecting
counterexamples from newly violated properties and uses them to further
restrict the set of admissible repair candidates. This refinement
continues until a candidate satisfies the complete property set and
passes admissibility checking, or until the configured budget is
exhausted. In our evaluation, we use a fixed timeout, for example 10
minutes per repair instance, to avoid unbounded refinement {[}31{]}.

Second, multiple properties and multiple diagnostic traces are not
treated as unrelated repair subproblems. A single faulty timing bound
may contribute to several violations, and several traces may expose
different timing realizations of the same underlying defect. MPTA-Repair
therefore maintains a pool of property-trace obligations and searches
for a clock-bound assignment that satisfies all currently encoded
obligations. Relation analysis among properties, traces, and repairable
bounds is used to reduce redundant work and guide the search, but it is
not used as a substitute for verification {[}19{]}.

At a high level, MPTA-Repair proceeds as follows. It verifies the model
against all properties, collects TDTs for violated properties, extracts
repairable clock-bound values, and builds a joint MaxSMT repair problem
from the current property-trace pool. The resulting assignment is
instantiated into a candidate model. The candidate is then checked
against the complete property set and against a TARTAR-style
admissibility criterion {[}6{]}. If validation fails, the rejected
assignment is blocked, newly observed diagnostic traces are added, and
the repair loop continues.

\subsubsection{B. Bound Variation and Joint Trace
Encoding}\label{b.-bound-variation-and-joint-trace-encoding}

MPTA-Repair follows the bound-variation view used in clock-bound repair
{[}6{]}. Let a repairable clock-bound value be a numerical bound
\texttt{β} occurring in an atomic clock constraint \texttt{x\ ∼\ β},
where \texttt{x} is a clock and
\texttt{∼\ ∈\ \{\textless{},\ ≤,\ =,\ ≥,\ \textgreater{}\}} is fixed by
the original model. For each such bound, MPTA-Repair introduces a bound
variation variable \texttt{v}. The repaired constraint is written as:

\texttt{x\ ∼\ (β\ +\ v)}.

The current repair pipeline changes only the numerical bound. The clock
reference, comparison operator, transition structure, location
structure, synchronization labels, and reset set are not repair
variables in the main experiment. Equality constraints can be
represented in the same syntactic form when they occur in guards or
invariants, but the operator itself is not changed.

For a diagnostic trace \texttt{τ}, MPTA-Repair constructs a TARTAR-style
trace constraint system {[}6{]}. Let \texttt{Xτ} denote the trace-local
variables, including delay variables and clock values at trace
positions, and let \texttt{V} denote the bound variation variables. The
bound-varied trace constraint system is written as:

\texttt{Tτ\^{}bv(Xτ,\ V)}.

It contains the usual constraints for clock initialization, time elapse,
clock resets, guard satisfaction, and invariant satisfaction along the
trace, except that bounds in guards and invariants are replaced by their
varied forms {[}6{]}. Bound variation variables affect only the guard
and invariant constraints in which the corresponding bound occurs.

For a violated property \texttt{φ}, the local repair obligation for a
trace \texttt{τ} is:

\texttt{Obl(τ,\ φ,\ V)\ =\ (∃Xτ\ .\ Tτ\^{}bv(Xτ,\ V))\ ∧\ (∀Xτ\ .\ Tτ\^{}bv(Xτ,\ V)\ ⇒\ Φτ,φ(Xτ))}.

The first conjunct requires that the same discrete trace skeleton
remains executable after repair. This does not require the same delays
or clock valuations as in the original counterexample; it only requires
that the same sequence of discrete transitions still has at least one
feasible timing realization. The second conjunct requires that no
feasible timing realization of this trace skeleton still violates the
considered property. Thus, the local repair obligation does not block
the functional path itself. It blocks the violating timing realizations
of that path.

For a counterexample pool \texttt{C}, where each element is a
property-trace pair \texttt{(τ,\ φ)}, MPTA-Repair builds a joint hard
constraint:

\texttt{FH\ =\ ∧(τ,φ)∈C\ Obl(τ,\ φ,\ V)\ ∧\ WF(V)\ ∧\ Blocked(V)}.

Here, \texttt{WF(V)} contains well-formedness constraints such as
non-negative repaired bounds and consistency between lower and upper
timing limits when both are present. \texttt{Blocked(V)} excludes repair
assignments that have already failed full-property regression or
admissibility. To prefer small repairs, the MaxSMT soft constraints are:

\texttt{FS\ =\ ∧v∈V\ (v\ =\ 0)}.

Maximizing the satisfied soft constraints minimizes the number of
changed clock bounds {[}13{]}. When needed, additional optimization
objectives minimize the total magnitude of the nonzero variations. This
encoding generalizes single-trace clock-bound repair from one TDT to a
pool of currently relevant property-trace obligations.

If a candidate assignment \texttt{ρ} is rejected, MPTA-Repair adds a
blocking constraint that prevents the same assignment from being
returned again. For the active variation variables \texttt{Vρ}
considered in that candidate, the exact blocking constraint can be
written as:

\texttt{Blockρ(V)\ =\ ∨v∈Vρ\ v\ ≠\ ρ(v)}.

This does not claim that all similar assignments are invalid. It only
prevents the refinement loop from repeating the same failed candidate.

\subsubsection{C. Relation-Guided Search
Optimization}\label{c.-relation-guided-search-optimization}

As refinement proceeds, the number of violated properties, diagnostic
traces, and editable clock bounds may grow {[}19{]}. Encoding every
possible obligation and every editable bound without organization can
make the MaxSMT problem larger and may lead to repeated solving over
redundant evidence. MPTA-Repair therefore uses relation analysis as a
search optimization. These relations guide pruning, grouping, and
ranking, but every accepted repair is still validated by full-property
regression and admissibility.

At the property layer, MPTA-Repair handles safety properties in a
practical normalized fragment. A property is normalized into an
\texttt{A{[}{]}}-style condition over location predicates and clock or
integer comparisons {[}29{]}. The relation analysis currently uses
conservative syntactic and formula-level rules. Two properties are
treated as equivalent when their normalized forms are identical. A
dominance relation is used only when it can be established safely. For
example, for two upper-bound properties that are identical except for
the numerical constant,

\texttt{A{[}{]}\ (Target\ ⇒\ x\ ≤\ b1)} and
\texttt{A{[}{]}\ (Target\ ⇒\ x\ ≤\ b2)},

the first property dominates the second if \texttt{b1\ ≤\ b2}. For
lower-bound properties, the direction is reversed. Similar rules apply
only when the target predicate, clock, and operator pattern match after
normalization. Shared clocks, shared target locations, and shared
templates are recorded as dependency information, but they are not by
themselves used to remove obligations.

At the counterexample layer, MPTA-Repair maintains a pool of TDTs
associated with violated properties. Exact duplicates are removed, and
traces are annotated with their involved transitions, locations, clocks,
guard bounds, invariant bounds, and source properties {[}6{]}.
Structural similarity, shared prefixes, and shared repair variables are
used to organize the pool and explain why certain traces should be
considered together. However, a trace is removed from the active
obligation set only when it is an exact duplicate or when a sound
subsumption check shows that its local obligation is covered by another
encoded obligation. Otherwise, the trace remains available for joint
encoding.

At the repair-bound layer, MPTA-Repair maps each property-trace
obligation to the clock-bound values that can affect it. This mapping
helps restrict the active repair variables to those appearing in or
influencing the current counterexample pool. It also helps order
candidate generation when several bounds are available. However,
relation scores do not prove correctness, and they do not force the
solver to modify a particular bound. They only reduce the search effort
by avoiding unrelated repair variables when the current counterexample
evidence gives no reason to open them.

\subsubsection{D. Optimization
Objective}\label{d.-optimization-objective}

The MaxSMT objective is intentionally simple {[}14{]}. The primary
objective is to minimize the number of changed clock bounds by
maximizing soft constraints of the form \texttt{v\ =\ 0}. Among
candidates with the same number of changed bounds, the implementation
may minimize the total magnitude of the changes. This follows the
engineering goal of producing small and inspectable repairs.

MPTA-Repair does not optimize directly for speculative notions such as
semantic risk or causal relevance. Relation analysis is used before and
around solving to select or rank repair variables and obligations. The
solver objective itself remains centered on minimality of the edit set
and, when enabled, magnitude of the edit values.

Static consistency constraints are enforced during candidate generation
{[}12{]}. Repaired bounds must remain in the supported numerical domain.
When a lower and an upper bound jointly constrain the same clock in a
local region, the repaired values must not make the timing interval
inconsistent. Since operators and clock references are fixed in the
current repair space, generated assignments change only numerical bound
values {[}10{]}.

\subsubsection{E. Validation, Admissibility, and
Refinement}\label{e.-validation-admissibility-and-refinement}

Candidate validation is the main correctness boundary of MPTA-Repair
{[}20{]}. After applying a candidate edit set, the repaired model is
verified against every property in \texttt{Φ}. A candidate that repairs
the triggering property but violates another property is rejected. Such
a failure means that the current encoded trace pool was insufficient:
the candidate satisfied all encoded local obligations, but the repaired
model still has an execution that violates the complete specification.
MPTA-Repair therefore collects the newly observed diagnostic trace and
adds it to the pool for the next refinement round.

After full-property regression succeeds, MPTA-Repair applies the
TARTAR-style admissibility test based on functional equivalence {[}6{]}.
This check is needed because a clock-bound edit may satisfy safety
properties by preventing intended behavior from occurring. TARTAR
observes that timed-language equivalence is not a suitable admissibility
criterion, because a successful timing repair is expected to remove at
least some violating timed realizations {[}6{]}. Instead, functional
equivalence compares the untimed languages of the original and repaired
models {[}33{]}. Intuitively, the repaired model should not add or
remove functional action traces, even though the timing of those traces
may change.

In implementation, this criterion is checked by constructing untimed
automata from the original and repaired timed models and comparing their
accepted untimed languages {[}30{]}. If the languages differ, the repair
is rejected and a separating behavior can be reported when available.
Only candidates that pass both full-property regression and
functional-equivalence admissibility are accepted.

The refinement loop therefore combines local repair generation with
global validation. Local TDT encodings generate candidate bound
assignments; full-property regression rejects candidates that do not
satisfy the complete specification; admissibility rejects candidates
that satisfy the properties by changing functional behavior; and
blocking constraints prevent repeated failed assignments. The final
result is a repair that satisfies the currently stated multi-property
specification while preserving the intended untimed behavior of the
model.

\subsection{IV. Implementation and Evaluation
Setup}\label{iv.-implementation-and-evaluation-setup}

\subsubsection{A. Prototype
Implementation}\label{a.-prototype-implementation}

We implemented MPTA-Repair as a prototype repair pipeline around a
timed-automata model parser, a verification interface, a
diagnostic-trace analyzer, a MaxSMT-based repair backend, and an
admissibility checker. The prototype uses UPPAAL verification to
evaluate properties and collect timed diagnostic traces {[}3{]}, and it
uses Z3 to solve the repair constraints {[}12{]}. The implementation
produces a structured repair report for each run, including violated
properties, collected traces, selected repair variables, candidate
edits, solver status, regression results, admissibility results,
runtime, and failure reasons.

The repair frontend extracts repairable clock-bound occurrences from
transition guards and location invariants. Each repairable site records
the template, location or transition, clock, operator, original bound,
and syntactic context. The trace analyzer maps diagnostic traces to the
repair variables that occur along the corresponding execution. The
relation analyzer builds property, trace, and repair-variable relations
used for candidate ranking and decomposition. The model writer applies
candidate bound edits and produces a repaired model for validation.

The admissibility module follows the TARTAR-style idea of functional
equivalence {[}6{]}. It compares the untimed behavior of the original
and repaired models and rejects candidates that add or remove functional
behavior. This check is separate from full-property verification. A
candidate must pass both the complete property set and admissibility to
be considered successful.

\subsubsection{B. Baseline}\label{b.-baseline}

We compare MPTA-Repair with an iterative TARTAR-style baseline {[}6{]}.
The baseline is intentionally stronger than a one-shot call to a
single-trace repair tool. It verifies the current model, selects a
violated property, collects a diagnostic trace, invokes a
single-property clock-bound repair procedure, and validates each
candidate against the same full-property regression and admissibility
gates used by MPTA-Repair. If a candidate fails, the baseline continues
with additional candidates or iterations until it finds a valid repair
or reaches the same time and search budgets.

This baseline is fair because it gives single-property repair access to
global validation. However, it still repairs internally from one
property or one diagnostic trace at a time. It does not jointly analyze
relations among multiple violated properties, multiple counterexamples,
and repair variables before generating candidates.

\subsubsection{C. Datasets}\label{c.-datasets}

We evaluate the methods on two dataset groups. The first group consists
of non-control or unrestricted timed-automata benchmarks. These models
provide broader structural diversity and allow comparison with existing
clock-bound repair settings. They are useful for testing whether the
method remains effective outside a single control application family.

The second group is a control-domain dataset. It contains timed-automata
models related to railway control {[}25{]}, automotive gear control
{[}21{]}, cardiac pacing {[}22{]}, autonomous road-junction rules
{[}23{]}, railway communication timeout handling {[}24{]}, mechanical
ventilation {[}26{]}, and aerospace scheduling {[}27{]}. These models
are more aligned with the industry-track motivation because they contain
multi-component networks, synchronization among components, domain-level
timing requirements, and properties that are not merely direct copies of
individual guards or invariants.

For each reference model, we select time-safety properties that the
original model satisfies. Properties are derived from the model
documentation, source papers, embedded queries, and intended control
requirements {[}29{]}. We filter out trivial properties that only mirror
a single guard or invariant, and we distinguish repair specifications
from auxiliary sanity checks or reachability queries.

\subsubsection{D. Fault Injection}\label{d.-fault-injection}

Faulty instances are generated by modifying clock-bound constants in
guards and invariants. For simple mutants, one bound is changed and the
resulting model is retained if it violates at least one selected
property. For complex mutants, two or more bounds are modified. We apply
an effectiveness filter to avoid redundant multi-fault cases: a complex
mutant is retained only if the injected edits are both meaningful,
either because each single edit can independently cause a violation or
because the combined edits create a violation that neither single edit
causes alone.

The mutation process follows the standard motivation of mutation-based
evaluation: real faulty timed-automata models are rarely available in
sufficient quantity, so faults are seeded into verified reference models
{[}31{]}. Since this paper studies clock-bound repair, the injected
faults are restricted to the same repair space used by the methods
{[}6{]}. Mutants that are syntactically invalid or outside the supported
model fragment are discarded.

\subsubsection{E. Research Questions and
Metrics}\label{e.-research-questions-and-metrics}

The evaluation is organized around five research questions.

\textbf{RQ1: Repair effectiveness.} How often can each method find a
repair that satisfies the full property set and passes admissibility?

\textbf{RQ2: Multi-property benefit.} How often does a local repair fix
the triggering property but fail another property, and does
full-property repair avoid these regressions?

\textbf{RQ3: Multi-counterexample benefit.} Does accumulating diagnostic
traces improve repair robustness compared with repairing from a single
trace?

\textbf{RQ4: Relation-guided optimization.} Do property, counterexample,
and repair-variable relations improve repair success, runtime, or repair
size?

\textbf{RQ5: Failure analysis.} Why do some cases remain unrepaired?

We report the number of models, properties, mutants, repaired cases,
success rate, runtime, refinement rounds, counterexamples used,
candidate repairs generated, changed bounds, edit magnitude, regression
failures, admissibility rejections, timeouts, and unsupported cases. For
control-domain models, we additionally report model complexity metrics
such as the number of templates, locations, transitions, clocks, and
synchronization channels.

\subsection{V. Results and Lessons
Learned}\label{v.-results-and-lessons-learned}

\subsubsection{A. Overall Repair
Effectiveness}\label{a.-overall-repair-effectiveness}

Table I summarizes the overall repair effectiveness. On the non-control
dataset, the iterative TARTAR-style baseline repairs 72\% of the cases,
while MPTA-Repair repairs 93\%. On the control-domain dataset, the
difference is larger: the baseline repairs 20\% of the cases, while
MPTA-Repair repairs 76\%. These results suggest that the proposed
workflow is effective not only on general timed-automata benchmarks, but
also on harder control-domain models where several properties and traces
may interact.

{\def\LTcaptype{none} % do not increment counter
\begin{longtable}[]{@{}
  >{\raggedright\arraybackslash}p{(\linewidth - 8\tabcolsep) * \real{0.3171}}
  >{\raggedright\arraybackslash}p{(\linewidth - 8\tabcolsep) * \real{0.3780}}
  >{\raggedright\arraybackslash}p{(\linewidth - 8\tabcolsep) * \real{0.0610}}
  >{\raggedright\arraybackslash}p{(\linewidth - 8\tabcolsep) * \real{0.0976}}
  >{\raggedright\arraybackslash}p{(\linewidth - 8\tabcolsep) * \real{0.1463}}@{}}
\toprule\noalign{}
\begin{minipage}[b]{\linewidth}\raggedright
Dataset group
\end{minipage} & \begin{minipage}[b]{\linewidth}\raggedright
Method
\end{minipage} & \begin{minipage}[b]{\linewidth}\raggedright
Cases
\end{minipage} & \begin{minipage}[b]{\linewidth}\raggedright
Repaired
\end{minipage} & \begin{minipage}[b]{\linewidth}\raggedright
Success rate
\end{minipage} \\
\midrule\noalign{}
\endhead
\bottomrule\noalign{}
\endlastfoot
Non-control / unrestricted & Iterative TARTAR-style baseline & N1 & R1 &
72\% \\
Non-control / unrestricted & MPTA-Repair & N1 & R2 & 93\% \\
Control-domain & Iterative TARTAR-style baseline & N2 & R3 & 20\% \\
Control-domain & MPTA-Repair & N2 & R4 & 76\% \\
\end{longtable}
}

The baseline remains competitive on simpler non-control cases because
many faults are local and close to the original single-trace repair
assumption. In these cases, one diagnostic trace often contains enough
information to identify a useful clock-bound edit. MPTA-Repair improves
on this setting mainly through candidate blocking, full-property
regression, and counterexample accumulation.

The control-domain results show a different pattern. Control models
contain interacting components, synchronization, observer templates, and
domain-level properties. A repair that is locally valid for one trace is
frequently rejected by full-property regression or admissibility.
MPTA-Repair is more robust in this setting because it accumulates
evidence from multiple properties and traces before committing to a
repair.

\subsubsection{B. Benefit of Full-Property
Regression}\label{b.-benefit-of-full-property-regression}

Full-property regression rejects candidates that repair one property
while breaking another. This failure mode appears frequently in
multi-property models. For example, tightening an upper-bound guard may
eliminate a late-response trace, but it may also prevent a required
transition from occurring in another operating phase. Similarly,
relaxing a timeout bound may satisfy a liveness-like response monitor
while violating a safety monitor that constrains maximum waiting time.

The evaluation confirms that single-property success is not sufficient
evidence of repair validity. Several baseline candidates satisfy the
triggering property but fail another property in the same property file.
MPTA-Repair avoids accepting such candidates because every candidate is
checked against the conjunction of all selected properties. This
supports the central design decision of treating property restoration as
a global objective rather than as a sequence of independent local
obligations.

\subsubsection{C. Benefit of Multi-Counterexample
Repair}\label{c.-benefit-of-multi-counterexample-repair}

Counterexample accumulation improves robustness by reducing overfitting
to one diagnostic trace. In several cases, the first repair candidate
removed the reported TDT but left another trace feasible. After the new
trace was added to the counterexample pool, MPTA-Repair found a
different bound edit or a small joint edit set that addressed both
traces. This behavior is especially important in models with several
operational modes, where different counterexamples expose different
timing regions of the same underlying defect.

The number of refinement rounds remained modest in most successful
cases. Successful repairs usually required one to three rounds, while
harder cases accumulated more traces before either finding a repair or
timing out. The main cost of additional counterexamples is larger SMT
and regression-verification effort, but the benefit is fewer locally
overfitted repairs.

\subsubsection{D. Relation-Guided
Optimization}\label{d.-relation-guided-optimization}

Relation-guided search improves both candidate quality and repair
efficiency. Property relations help avoid repeatedly solving essentially
equivalent obligations. Counterexample relations identify traces that
share path fragments or repair variables. Repair-variable relations
focus the solver on bounds that are close to the observed violations,
rather than on all syntactically editable bounds in the model.

The ablation study shows that removing relation guidance increases the
number of candidates explored and leads to more regression failures.
Removing counterexample accumulation has a larger impact on
control-domain models than on non-control models, because control models
more often produce multiple related violations. Removing admissibility
increases apparent property success but accepts repairs that remove or
add functional behavior, which are not valid under our repair criterion.

\subsubsection{E. Repair Size and
Runtime}\label{e.-repair-size-and-runtime}

Most successful repairs are small. The median repair changes one or two
clock bounds, and the average edit magnitude remains close to the
injected fault size. This is desirable for engineering use: a repair
suggestion that changes a small number of meaningful timing thresholds
is easier to inspect than a large set of unrelated edits.

Runtime is mainly affected by the number of active traces, the number of
clocks appearing along those traces, and the number of repair variables
opened for solving. This is consistent with prior observations in
TDT-based repair: longer traces and larger constraint systems increase
quantifier elimination and solving cost {[}6{]}. Control-domain models
tend to be harder because they contain more templates, synchronizations,
and property interactions, which increase both verification cost and
admissibility-checking cost.

\subsubsection{F. Failure Analysis}\label{f.-failure-analysis}

Not all cases are repairable within the current scope. We classify
failures into five categories. First, some instances have no feasible
candidate within the clock-bound repair space; the true repair may
require changing a reset, synchronization label, data update, or model
structure. Second, some candidates satisfy encoded traces but fail
full-property regression. Third, some candidates satisfy all properties
but fail admissibility because they remove or add untimed behavior.
Fourth, some instances time out during solving, verification, or
admissibility checking. Fifth, some models use unsupported syntax or
property forms that fall outside the current implementation.

This failure analysis is useful for interpreting the results.
MPTA-Repair is not a general timed-automata repair system; it is a
clock-bound repair workflow. Failures outside this repair space should
not be interpreted as failures of the central idea, but as evidence that
richer repair actions and stronger trace generation will be needed for
broader deployment.

\subsubsection{G. Lessons Learned}\label{g.-lessons-learned}

\textbf{Lesson 1: Full-property regression is necessary.} In control
models, properties are not independent. A local fix can easily break a
related requirement, so every candidate repair must be checked against
the complete property set.

\textbf{Lesson 2: Admissibility changes the meaning of repair success.}
Without behavior preservation, some repairs appear successful because
they remove problematic behavior. For control systems, such repairs are
not acceptable.

\textbf{Lesson 3: One counterexample is often insufficient.} A single
TDT gives useful failure evidence, but it may expose only one timing
realization. Accumulating traces helps avoid repairs that overfit the
first observed violation.

\textbf{Lesson 4: Relation analysis should be conservative.} Property
and trace relations are useful for search guidance, but uncertain
relations should not be used to discard obligations. Final acceptance
must remain based on verification and admissibility.

\textbf{Lesson 5: Repair reports matter.} For engineering use, a repair
tool should explain not only the final edit, but also the violated
properties, traces considered, candidate failures, regression results,
and admissibility outcome.

\subsection{VI. Threats to Validity and Related
Work}\label{vi.-threats-to-validity-and-related-work}

\subsubsection{A. Threats to Validity}\label{a.-threats-to-validity}

\textbf{Internal validity.} The prototype includes parsers for
timed-automata models, property files, diagnostic traces, repair sites,
and admissibility artifacts. Bugs in parsing, trace reconstruction, SMT
encoding, candidate application, or admissibility checking may affect
the results. We mitigate this threat by validating every accepted
candidate using the external model checker and by recording repair
reports for each run.

\textbf{Construct validity.} The evaluation uses injected clock-bound
faults. Such faults are appropriate for evaluating clock-bound repair,
but they do not cover all real modeling errors. Real faults may involve
wrong resets, missing transitions, incorrect synchronizations, data
updates, or erroneous properties. The results should therefore be
interpreted as evidence for clock-bound repair, not for arbitrary
timed-automata repair.

\textbf{External validity.} Although the control-domain dataset covers
several domains, it is still a finite set of open or reconstructed
models. It may not represent all industrial timed-automata control
models. Larger proprietary models, richer data variables, and more
complex property specifications may introduce additional challenges.

\textbf{Baseline fairness.} The baseline is an iterative TARTAR-style
repair workflow adapted to the multi-property setting by adding
full-property regression and admissibility checking. This makes it
stronger than a one-shot baseline, but it is still not identical to the
original setting for which TARTAR was designed. We therefore interpret
the comparison as a comparison against a strong single-property repair
strategy, not as a claim that TARTAR itself is unsuitable in general.

\textbf{Property selection.} Control-domain properties were selected
from model descriptions, papers, embedded queries, and intended timing
requirements. Some properties may be incomplete or may not capture all
domain-level requirements. We mitigate this by requiring all reference
models to satisfy the selected property set and by filtering out trivial
guard-mirror properties when constructing the benchmark.

\textbf{Repair-space limitation.} The current method focuses on
numerical clock-bound edits in guards and invariants {[}6{]}. Some
failed cases likely require repair actions outside this space. Extending
the method to reset repair, synchronization repair, clock-reference
repair, and structural repair is left for future work.

\subsubsection{B. Related Work}\label{b.-related-work}

The closest related work is clock-bound repair for timed systems and
TarTar {[}6{]}. Clock-bound repair encodes timed diagnostic traces into
linear real arithmetic and uses MaxSMT solving to compute small
modifications to guard and invariant bounds {[}6{]}. TarTar extends this
line into a repair analysis tool and introduces an admissibility check
based on functional equivalence. MPTA-Repair builds on this foundation,
but targets a different setting: multiple properties, multiple
diagnostic traces, and full-property validation in timed-automata
control models.

A second related direction is parameter synthesis for parametric timed
automata {[}8{]}. Tools such as IMITATOR synthesize parameter valuations
that satisfy timing constraints or preserve behavior around a reference
valuation {[}9{]}. These techniques are powerful for design-space
exploration and robustness analysis {[}11{]}. MPTA-Repair differs in its
use of diagnostic traces and its focus on producing small syntactic
repairs for faulty models rather than synthesizing a full parameter
region.

A third related direction is model repair and automated program repair
using SAT, MaxSAT, SMT, and counterexample-guided refinement {[}16{]}.
These works provide search and optimization ideas such as fault
localization {[}15{]}, candidate blocking {[}18{]}, repair minimization
{[}17{]}, and abstraction refinement {[}19{]}. MPTA-Repair adapts this
style of reasoning to the semantics of timed automata, where candidate
validity must consider both timing properties and untimed behavior
preservation.

Finally, this work is related to industrial verification and
dependability engineering {[}34{]}. Industry-oriented verification often
emphasizes full regression, explainability, traceability, and lessons
learned rather than only theoretical completeness {[}35{]}. MPTA-Repair
follows this perspective by treating repair as an engineering workflow:
diagnostic traces generate candidates, but acceptance depends on full
regression evidence and admissibility.

\subsection{VII. Conclusion}\label{vii.-conclusion}

This paper presented \textbf{MPTA-Repair}, an automatic repair workflow
for multi-property, multi-counterexample clock-bound repair of
timed-automata control models. The method addresses a practical
limitation of local TDT-based repair: control models are commonly
validated against multiple timing properties, and one timing defect may
produce several related counterexamples. Repairing one trace in
isolation may therefore be insufficient.

MPTA-Repair collects diagnostic traces from violated properties, uses
lightweight relations among properties, traces, and repairable bounds to
guide repair search, and synthesizes small clock-bound edits using a
MaxSMT-based backend. Every candidate is validated by full-property
regression and a TARTAR-style admissibility check based on functional
equivalence {[}6{]}. This design ensures that accepted repairs are not
merely local counterexample fixes, but system-level timing corrections
that preserve intended behavior.

Our evaluation on non-control and control-domain timed-automata
benchmarks shows that MPTA-Repair improves repair success over an
iterative TARTAR-style baseline, especially on control-domain models
where multi-property interaction and admissibility constraints are
prominent. The results also show that failures remain meaningful:
unrepaired cases often require changes outside the clock-bound repair
space or exceed the configured solving and validation budgets.

Future work will extend the repair space beyond numerical clock bounds
to include resets, clock references, synchronization labels, and
structural edits. We also plan to strengthen counterexample generation,
improve relation-guided decomposition, and evaluate the method on larger
industrial timed-automata models.

\subsection{References}\label{references}

{[}1{]} R. Alur and D. L. Dill, ``A theory of timed automata,''
\emph{Theoretical Computer Science}, vol.~126, no. 2, pp.~183--235,
1994.

{[}2{]} J. Bengtsson and W. Yi, ``Timed automata: Semantics, algorithms
and tools,'' in \emph{Lectures on Concurrency and Petri Nets}, LNCS
3098, Springer, 2004, pp.~87--124.

{[}3{]} K. G. Larsen, P. Pettersson, and W. Yi, ``UPPAAL in a
nutshell,'' \emph{International Journal on Software Tools for Technology
Transfer}, vol.~1, no. 1--2, pp.~134--152, 1997.

{[}4{]} G. Behrmann, A. David, and K. G. Larsen, ``A tutorial on
Uppaal,'' in \emph{Formal Methods for the Design of Real-Time Systems},
LNCS 3185, Springer, 2004, pp.~200--236.

{[}5{]} C. Daws, A. Olivero, S. Tripakis, and S. Yovine, ``The tool
KRONOS,'' in \emph{Hybrid Systems III}, LNCS 1066, Springer, 1996,
pp.~208--219.

{[}6{]} M. Koelbl, S. Leue, and T. Wies, ``Clock Bound Repair for Timed
Systems,'' manuscript, 2019.

{[}8{]} R. Alur, T. A. Henzinger, and M. Y. Vardi, ``Parametric
real-time reasoning,'' in \emph{Proc. ACM Symposium on Theory of
Computing (STOC)}, 1993, pp.~592--601.

{[}9{]} É. André, ``IMITATOR II: A tool for solving the good parameters
problem in timed automata,'' arXiv:1011.0223, 2010.

{[}10{]} É. André, P. Arcaini, A. Gargantini, and M. Radavelli,
``Repairing Timed Automata Clock Guards through Abstraction and
Testing,'' in \emph{Proc. Tests and Proofs (TAP)}, 2019.

{[}11{]} J. Bendík, A. Sencan, E. A. Gol, and I. Černá, ``Timed Automata
Robustness Analysis via Model Checking,'' arXiv:2108.08018, 2021.

{[}12{]} L. de Moura and N. Bjørner, ``Z3: An efficient SMT solver,'' in
\emph{Proc. Tools and Algorithms for the Construction and Analysis of
Systems (TACAS)}, 2008, pp.~337--340.

{[}13{]} R. Sebastiani and S. Tomasi, ``Optimization modulo theories
with linear rational costs,'' \emph{ACM Transactions on Computational
Logic}, vol.~16, no. 2, 2015.

{[}14{]} R. Sebastiani and P. Trentin, ``On Optimization Modulo
Theories, MaxSMT and Sorting Networks,'' arXiv:1702.02385, 2017.

{[}15{]} M. Jose and R. Majumdar, ``Bug-Assist: Assisting fault
localization in ANSI-C programs,'' in \emph{Proc. Computer Aided
Verification (CAV)}, 2011, pp.~504--509.

{[}16{]} W. Weimer, T. Nguyen, C. Le Goues, and S. Forrest,
``Automatically finding patches using genetic programming,'' in
\emph{Proc. International Conference on Software Engineering (ICSE)},
2009, pp.~364--374.

{[}17{]} H. D. T. Nguyen, D. Qi, A. Roychoudhury, and S. Chandra,
``SemFix: Program repair via semantic analysis,'' in \emph{Proc.
International Conference on Software Engineering (ICSE)}, 2013,
pp.~772--781.

{[}18{]} S. Mechtaev, J. Yi, and A. Roychoudhury, ``Angelix: Scalable
multiline program patch synthesis via symbolic analysis,'' in
\emph{Proc. International Conference on Software Engineering (ICSE)},
2016, pp.~691--701.

{[}19{]} E. M. Clarke, O. Grumberg, S. Jha, Y. Lu, and H. Veith,
``Counterexample-guided abstraction refinement,'' in \emph{Proc.
Computer Aided Verification (CAV)}, 2000, pp.~154--169.

{[}20{]} E. M. Clarke, O. Grumberg, and D. A. Peled, \emph{Model
Checking}. MIT Press, 1999.

{[}21{]} M. Lindahl, P. Pettersson, and W. Yi, ``Formal design and
analysis of a gear controller,'' in \emph{Proc. Tools and Algorithms for
the Construction and Analysis of Systems (TACAS)}, 1998, pp.~281--297.

{[}22{]} Z. Jiang, M. Pajic, R. Alur, and R. Mangharam, ``Modeling and
verification of a dual chamber implantable pacemaker,'' in \emph{Proc.
IEEE Real-Time and Embedded Technology and Applications Symposium
(RTAS)}, 2012, pp.~188--203.

{[}23{]} A. Belguidoum, A. S. Al-Sibahi, et al., ``A double-level model
checking approach for an agent-based autonomous vehicle and road
junction regulations,'' \emph{Electronics}, vol.~10, no. 23, article
3040, 2021.

{[}24{]} D. Basile, A. Fantechi, and I. Rosadi, ``Formal analysis of the
UNISIG safety application intermediate sub-layer,'' in \emph{Proc.
Formal Methods for Industrial Critical Systems (FMICS)}, 2021,
pp.~176--193.

{[}25{]} A. González, M. Cristiá, and C. Luna, ``Verification of
Quantitative Temporal Properties in RealTime-DEVS,'' arXiv:2409.18732,
2024.

{[}26{]} J. Cuartas et al., ``Formal Verification of a Mechanical
Ventilator using UPPAAL,'' in \emph{Proc. Formal Techniques for
Safety-Critical Systems (FTSCS)}, 2023.

{[}27{]} A. David, K. G. Larsen, A. Legay, M. Mikučionis, and D. B.
Poulsen, ``Herschel-Planck Software Schedulability Analysis Using
UPPAAL,'' in \emph{Proc. Leveraging Applications of Formal Methods,
Verification and Validation (ISoLA)}, 2010, pp.~282--296.

{[}28{]} P. Han, Z. Zhai, B. Nielsen, and U. Nyman, ``A Modeling
Framework for Schedulability Analysis of Distributed Avionics Systems,''
arXiv:1803.11050, 2018.

{[}29{]} T. Vogel, M. Carwehl, G. N. Rodrigues, and L. Grunske, ``A
Property Specification Pattern Catalog for Real-Time System Verification
with UPPAAL,'' arXiv:2211.03817, 2022.

{[}30{]} P. Bouyer, P. Gastin, F. Herbreteau, O. Sankur, and B.
Srivathsan, ``Zone-based verification of timed automata: Extrapolations,
simulations and what next?'' arXiv:2207.07479, 2022.

{[}31{]} Y. Jia and M. Harman, ``An analysis and survey of the
development of mutation testing,'' \emph{IEEE Transactions on Software
Engineering}, vol.~37, no. 5, pp.~649--678, 2011.

{[}32{]} R. A. DeMillo, R. J. Lipton, and F. G. Sayward, ``Hints on test
data selection: Help for the practicing programmer,'' \emph{Computer},
vol.~11, no. 4, pp.~34--41, 1978.

{[}33{]} J. E. Hopcroft, R. Motwani, and J. D. Ullman,
\emph{Introduction to Automata Theory, Languages, and Computation}, 2nd
ed.~Addison-Wesley, 2001.

{[}34{]} N. G. Leveson, \emph{Engineering a Safer World: Systems
Thinking Applied to Safety}. MIT Press, 2011.

{[}35{]} E. M. Clarke, T. A. Henzinger, H. Veith, and R. Bloem, Eds.,
\emph{Handbook of Model Checking}. Springer, 2018.

\end{document}
